\section{一ブロックの相対量子化とPBW性}
\label{sec:pg-quantization-foundations}
ここから代数的量子化を構成する。周囲の二次代数を$A$、
その正則中心三次商を$B=A/(\Omega)$と書く。
まず一ブロックの全次数PBW性と中心正則拡大の入力を準備する。
中心公式の有限次数の検算、全次数のPBW性、および
中心元の左右正則性は、それぞれ別の主張として証明する。

\subsection{次数付き量子化とAS条件}

\begin{definition}
連結次数付き可換代数 $C$ に次数零の Poisson 括弧があるとする。
$0\in U\subset\mathbb A^1_\hbar$、$K=\OO(U)$ 上の次数付き量子化とは、
$\deg\hbar=0$、$\mathcal C_0=K$、各 $\mathcal C_d$ が有限自由な
次数付き代数 $\mathcal C$ と $\mathcal C/(\hbar)\simeq C$ であって、
\begin{equation}
\frac{[\widetilde a,\widetilde b]}\hbar\bmod\hbar=\{a,b\}
\end{equation}
を満たすものをいう。特殊 fiber の可換性と平坦性により持上げの選択に依存しない。
\end{definition}
連結次数付き代数 $A$ が Koszul であるとは、自明加群 $\CC=A/A_{\geq1}$ が線形自由分解をもつことをいう。
左右の injective dimension が $d$ で、左右それぞれ
\begin{equation}
\Ext_A^i(\CC,A)=
\begin{cases}\CC(\ell)&i=d,\\0&i\neq d\end{cases}
\end{equation}
ならば、次元 $d$、AS parameter $\ell$ の AS-Gorenstein 代数という。
さらに大域次元 $d$ と有限 GK 次元をもてば AS-regular という。
次数移動は $M(\ell)_j=M_{j+\ell}$ とする。
以下の Ext は全内部次数を含み、injective/global dimension は通常の左右加群圏で取る。

\subsection{ブロック代数と PBW 性}
$R=\CC[u_1,\ldots,u_m]$ とし、$\sym_R$ は中心係数を外へ出して active 変数の語を平均する対称化である。
一ブロック $f(x_1,x_2,x_3,u)$ に対し
\begin{equation}
A_f=R[\hbar]\langle x_1,x_2,x_3\rangle/
\bigl([x_2,x_3]-\hbar\sym_R(f_{x_1}),\ [x_3,x_1]-\hbar\sym_R(f_{x_2}),\
[x_1,x_2]-\hbar\sym_R(f_{x_3})\bigr).\label{eq:pg-block-algebra}
\end{equation}
これは相対 cyclic potential
\begin{equation}
W=x_1x_2x_3-x_2x_1x_3-\hbar\sym_R(f)\label{eq:pg-potential}
\end{equation}
の active cyclic derivatives で定まる。
cyclic word は自由代数を交換子の線形包で割った類であり、
$\partial_i^{\mathrm{cyc}}(a_1\cdots a_r)
=\sum_{a_j=x_i}a_{j+1}\cdots a_ra_1\cdots a_{j-1}$ とする。

\begin{proposition}\label{prop:pg-PBW}
有限個のブロックに共通な $0\in U_{\mathrm{reg}}\subset\mathbb A^1$ を選び、以下を満たせる。
$g_s(x_s)=f_s(x_s,0)$ とおくと、全 $\lambda\in U_{\mathrm{reg}}$ で
\begin{equation}
S_{s,\lambda}=\CC\langle x_{s1},x_{s2},x_{s3}\rangle/
([x_{s2},x_{s3}]-\lambda\sym((g_s)_{x_{s1}}),\text{循環置換})
\end{equation}
は三次元 Koszul AS-regular domain であり、3-Calabi--Yau である。
その Hilbert 級数は $(1-t)^{-3}$。
active 次数で filter し、$u_\alpha$ の filtration degree を零にすると
\begin{equation}
\operatorname{gr}_xA_{f_s,\lambda}\simeq S_{s,\lambda}\otimes_\CC R.
\end{equation}
族全体にも対応する同型があり、$A_{f_s}|_{U_{\mathrm{reg}}}$ は
$R_{\mathrm{reg}}=\OO(U_{\mathrm{reg}})[u]$ 上自由である。
\end{proposition}
\begin{proof}
一つの $g$ について関係式の列は $M_\lambda x$ で表される。ただし
\begin{equation}
M_\lambda=
\begin{pmatrix}
-\frac\lambda2g_{11}&-x_3-\frac\lambda2g_{12}&x_2-\frac\lambda2g_{13}\\
x_3-\frac\lambda2g_{21}&-\frac\lambda2g_{22}&-x_1-\frac\lambda2g_{23}\\
-x_2-\frac\lambda2g_{31}&x_1-\frac\lambda2g_{32}&-\frac\lambda2g_{33}
\end{pmatrix}.
\end{equation}
反対称成分が三つの交換子を与えるので関係式は独立である。
potential 由来の表示は\cite[Corollary 2.7]{PMoriSmith}の意味で standard であり、
同論文 Theorem 1.4に再掲されたArtin--Tate--Van den Berghの判定により、AS-regular 性は全 $\PP^2$ 上で $\rank M_\lambda\geq2$ と同値である。
階数一以下の incidence は $\mathbb A^1\times\PP^2$ の閉集合で、その底への像も閉じている。
$M_0$ は全 $\PP^2$ 上で階数二の交代行列なので、像を除けば原点を含む開集合を得る。
有限個のブロックについて交わりを取る。
3-CY 性と domain 性は\cite[Corollary 2.9]{PMoriSmith}による。
三生成元三二次関係式の三次元 AS-regular 代数では、最小分解を双対化すると最後の shift は三となり、
\begin{equation}
0\longrightarrow S_{s,\lambda}(-3)\longrightarrow S_{s,\lambda}(-2)^3
\longrightarrow S_{s,\lambda}(-1)^3\longrightarrow S_{s,\lambda}
\longrightarrow\CC\longrightarrow0.
\end{equation}
これが Koszul 性と Hilbert 級数を与える。

$\lambda$ を固定する。$u=a$ への特殊化では、\eqref{eq:pg-potential}は
$g$ の三次 potential に active 次数二以下の項を加えたものである。
\cite[Theorem 1.1(i)]{PBergerTaillefer}により、その代数は $S_{s,\lambda}$ の PBW 変形である。
$S_{s,\lambda}$ の斉次基底を語の線形結合として持ち上げる。
最高 active 次数の関係式から次数帰納法で、これらは $A_{f_s,\lambda}$ を $R$ 上生成する。
線形関係があれば全 $u=a$ で PBW 基底間の関係となり、係数多項式が全点で零だから元の関係も零となる。
従って所要の PBW 性と自由性を得る。

族の斉次 $S_{s,U_{\mathrm{reg}}}$ の各次数部分は底上有限生成で、全 fiber の次元が一定である。
底環は $\CC[\hbar]$ の局所化で PID なので、各部分は有限自由である。
torsion があればその台で fiber 次元が増えるからである。
その各次数の基底を自由代数へ持ち上げ、同じ帰納法と全 $(\lambda,a)$ への特殊化を使えば、族上の自由性も従う。
\end{proof}
\begin{remark}\label{rem:pg-regular-all}
$g_s$ が滑らかな平面三次曲線を定めるとき、
\cite[Theorem 4.2(1)]{PMoriSmith}により $S_{s,\lambda}$ の AS-regular 性のために除く $\lambda$ はない。
$g_s=0$ の場合も常に多項式環である。
中心元の代数化・正則性のためには、以下でさらに開集合を縮める。
ambient の AS-regular 性は、その中心超曲面商の圏論的滑らかさとは別条件である。
\end{remark}
以下
$A_{\mathrm{reg}}=A_{f_1}\otimes_{R_{\mathrm{reg}}}\cdots\otimes_{R_{\mathrm{reg}}}A_{f_k}$
とする。異なるブロックは可換で、各全次数部分は底上有限自由であり、
$A_{\mathrm{reg}}/(\hbar)\simeq S$。関係式の一次項は\eqref{eq:pg-bracket}である。

\subsection{代数的中心元の規格化と正則中心拡大}
任意の滑らかな三次ブロックにおける中心元の存在とCramer表示は、
後の定理\ref{thm:pg-universal-determinant}で証明する。
同定理が使う古典的Casimir核の計算は、
三変数制限が特異・零の場合も成立する。
さらに複数ブロックを持つ滑らかな$F$では
補題\ref{lem:pg-central-splitting-smooth}によって
中心変数だけの三次項を配分し、各ブロックを
\emph{全変数に関して滑らかな三次式}に選べる。
これらを合わせれば全Poisson構造に対して
有限次数の行列式演算で代数的中心三次元を得る。

\begin{lemma}[中心三次元の偶性]\label{lem:pg-center-parity}
命題\ref{prop:pg-PBW}の代数的族$A_f$について、
$\hbar\mapsto-\hbar$で不変な$0$を含む開集合へ必要なら縮める。
生成元に
\begin{equation}
 x_i^\dagger=x_i,\qquad u_\alpha^\dagger=u_\alpha,\qquad
 \hbar^\dagger=-\hbar,\qquad(ab)^\dagger=b^\dagger a^\dagger
\end{equation}
を指定すると、$\CC$-線形の反自己同型$\dagger$が定まる。
中心元$\Omega\in Z(A_f)_3$、$\Omega\bmod\hbar=f$があれば、
同じ剰余をもつ中心元を$\Omega^\dagger=\Omega$となるように選べる。
特に形式的中心元も同様に選べ、Weyl対称化表示では
\begin{equation}
 \widehat\Omega
 =\sym_R\bigl(f+\hbar^2 f_2+\hbar^4 f_4+\cdots\bigr),
 \qquad f_{2r}\in R[x_1,x_2,x_3]_3
\end{equation}
である。ここで添字$3$は中心変数$u_\alpha$も次数一と
数えた全次数を意味する。
\end{lemma}
\begin{proof}
定義関係式
$r_i=[x_j,x_k]-\hbar\sym_R(f_{x_i})$に対し、
$r_i^\dagger=-r_i$である。したがって反自己同型が商へ降り、
中心を保つ。$\dagger^2=1$なので
$(\Omega+\Omega^\dagger)/2$は中心であり、$\hbar=0$では$f$に一致する。
完備化でも同じ操作が可能である。

全次数三のWeyl対称化
$\sym_R:R[x_1,x_2,x_3]_3[[\hbar]]\to\widehat A_{f,3}$
は$\hbar=0$で恒等写像となり、
PBW性と完備Nakayamaにより同型である。
代数的族についてもこの有限自由加群間の写像は
原点で可逆なので、さらに開集合を縮めれば同型となる。
対称化した単語は反転で不変であるから、
$\dagger$はその係数の$\hbar\mapsto-\hbar$に対応する。
したがって不変な中心元の対称化係数は$\hbar$の偶冪だけを含む。
\end{proof}

\begin{proposition}[普遍的なCasimirの二次補正]\label{prop:pg-universal-casimir-jet}
$f\in R[x_1,x_2,x_3]$を全次数三の任意の多項式とし、
\begin{equation}
 H_f=\det\left(\frac{\partial^2f}
 {\partial x_i\partial x_j}\right)_{i,j=1}^3,\qquad
 \Theta_f=\sym_R\left(f-\frac{\hbar^2}{12}H_f\right)\in A_f
\end{equation}
とおく。命題\ref{prop:pg-PBW}の開集合上で
\begin{equation}
 [\Theta_f,x_i]\in\hbar^5 A_f\qquad(i=1,2,3),
 \qquad [\Theta_f,u_\alpha]=0.
 \label{eq:pg-casimir-jet}
\end{equation}
従って$\Theta_f$は$\hbar^5$を法として中心元であり、
法$\hbar$で$f$となる。
さらに補題\ref{lem:pg-center-parity}で正規化した
任意の形式的中心元について、$\hbar^2$係数$f_2$は
\begin{equation}
 \left\{f_2+\frac1{12}H_f,x_i\right\}_f=0\quad(i=1,2,3)
 \label{eq:pg-jet-casimir-condition}
\end{equation}
を満たす。
最後の条件は$f_2=-H_f/12$が
\emph{全次数で厳密な中心元に延長する}ことを主張しない。
\end{proposition}
\begin{proof}
まず$\hbar^4$を法とした恒等式を示す。
係数が代数的に独立な二十個の不定元である
\begin{equation}
 f_{\mathrm{univ}}=\sum_{a+b+c\le3}
 t_{abc}x_1^a x_2^b x_3^c
\end{equation}
を考える。これは中心係数環を任意に特殊化できる普遍式である。
自由語では、隣り合った逆順の文字を
\begin{equation}
 x_2x_1\rightsquigarrow x_1x_2-\hbar\sym(f_{x_3}),\quad
 x_3x_1\rightsquigarrow x_1x_3+\hbar\sym(f_{x_2}),\quad
 x_3x_2\rightsquigarrow x_2x_3-\hbar\sym(f_{x_1})
 \label{eq:pg-jet-reductions}
\end{equation}
で置き換える。各置換は代数の定義関係式そのものである。
$\hbar$次数が三以下の範囲では、
逆転数が減るか$\hbar$次数が増えるため、還元は必ず停止する。
$\sym(f_{\mathrm{univ}}-\hbar^2H_{f_{\mathrm{univ}}}/12)$と
各$x_i$との交換子を展開して\eqref{eq:pg-jet-reductions}で還元すると、
次数三以下の各係数が$\mathbb Q[t_{abc}]$内で零になる。
これを古典Jacobian括弧$\{-,-\}_f$で書き直すと、
より明示的な普遍恒等式
\begin{equation}
 [\sym_R(f),x_i]\equiv
 \frac{\hbar^3}{12}\sym_R(\{H_f,x_i\}_f)
 \pmod{\hbar^4 A_f}
 \label{eq:pg-jet-hessian-identity}
\end{equation}
を得る。実際、$\sym_R(H_f)$の$x_i$との交換子の一次項は
$\hbar\sym_R(\{H_f,x_i\}_f)$なので、
この恒等式は$\Theta_f$の可換性の計算と同値である。
二十個の独立係数を保持して全係数を厳密に展開する
再現可能な検算コードは
\texttt{computations/verify\_cubic\_casimir\_jet.py}である。
数値の特殊化や標本点での検査ではない。
よって任意の中心係数環への特殊化で
$[\Theta_f,x_i]\in\hbar^4A_f$となる。

一方、$\Theta_f^\dagger=\Theta_f$であり、
$[\Theta_f,x_i]^\dagger=-[\Theta_f,x_i]$である。
PBW性により$\hbar$は非零因子なので、
$[\Theta_f,x_i]=\hbar^4b_i$と書くと
$b_i^\dagger=-b_i$となる。
法$\hbar$では反転は可換多項式環の恒等写像だから
$b_i\bmod\hbar=0$、従って$b_i\in\hbar A_f$。
これで法$\hbar^5$の結論を得る。
中心変数との可換性は定義から従う。

最後に$\widehat\Omega=\sym_R(f+\hbar^2f_2+O(\hbar^4))$
との交換子を比較する。
前二段の恒等式および量子化の一次Poisson条件から
$\hbar^3\{f_2+H_f/12,x_i\}_f=0\bmod\hbar^4$
となり、PBW性から\eqref{eq:pg-jet-casimir-condition}が従う。
\end{proof}
\begin{remark}[Hesse公式との照合]
$f=(x_1^3+x_2^3+x_3^3)/3+t x_1x_2x_3$なら
\begin{equation}
 H_f=2(t^3+4)x_1x_2x_3-2t^2(x_1^3+x_2^3+x_3^3).
\end{equation}
したがって$-\hbar^2H_f/12$は
定理\ref{thm:explicit-hesse-center}の
中心式の$\hbar^2$係数と一致する。
これは既知のSklyanin公式との整合性チェックであって、
一般$f$の全高次中心公式が同じ形になるという主張ではない。
\end{remark}

\begin{lemma}\label{lem:pg-central-linear}
$D$ を連結次数付き代数、$z\in D_1$ を中心正則元とする。
$C=D/(z)$ が次元 $d$、AS parameter $d$ の Koszul AS-regular なら、
$D$ は次元 $d+1$、AS parameter $d+1$ の Koszul AS-regular である。
\end{lemma}
\begin{proof}
$0\to D(-1)\xrightarrow{z}D\to C\to0$ より、
$\operatorname{Tor}^D_q(C,\CC)$ は $q=0,1$ でそれぞれ $\CC,\CC(-1)$、他は零である。
change-of-rings spectral sequence
\begin{equation}
E^2_{p,q}=\operatorname{Tor}^C_p(\operatorname{Tor}^D_q(C,\CC),\CC)
\Longrightarrow\operatorname{Tor}^D_{p+q}(\CC,\CC)
\end{equation}
の各項は内部次数 $p+q$ に集中する。
微分は内部次数を保ち全ホモロジー次数を一減らすため全て零であり、
$D$ は Koszul、$j>d+1$ で $\operatorname{Tor}^D_j(\CC,\CC)=0$、$j=d+1$ では非零である。
$D^!$ を二次双対とすると、Koszul 双加群複体の項は $D\otimes(D_i^!)^*\otimes D$。
全内部次数を固定し右の $D$ 次数で filter すれば有限 filtration となり、
随伴複体の正ホモロジーは左 Koszul 分解で消える。
よって増大双加群複体は長さ $d+1$ の完全列であり、右加群として分裂完全なので
任意の左加群とテンソルしても完全である。従って通常の大域次元も $d+1$。
さらに $\RHom_D(C,D)\simeq C(1)[-1]$ と導来随伴から
\begin{equation}
\Ext_D^{i+1}(\CC,D)\simeq\Ext_C^i(\CC,C)(1).
\end{equation}
右加群でも同様で、左右 AS 条件と parameter を得る。
$\operatorname{Hilb}_D=\operatorname{Hilb}_C/(1-t)$ なので GK 次元も有限である。
\end{proof}
\begin{lemma}\label{lem:pg-central-cubic}
$A$ を次元 $N$、AS parameter $N$ の AS-regular 代数、$\eta\in A_3$ を中心正則元とする。
$B=A/(\eta)$ は左右 injective dimension $N-1$、AS parameter $N-3$ の AS-Gorenstein 代数で、
$\operatorname{Hilb}_B=(1-t^3)\operatorname{Hilb}_A$ である。
\end{lemma}
\begin{proof}
$0\to A(-3)\xrightarrow{\eta}A\to B\to0$ が Hilbert 級数を与える。
導来随伴により任意の左 $B$-加群 $M$ について
$\Ext_B^i(M,B)\simeq\Ext_A^{i+1}(M,A)(-3)$。
右辺は $i\geq N$ で零、$M=\CC$、$i=N-1$ で $\CC(N-3)$ となる。
よって injective dimension はちょうど $N-1$ で AS 条件も成立する。右側も同様。
\end{proof}

